\documentclass[11pt, oneside]{book}   	% use "amsart" instead of "article" for AMSLaTeX format
\usepackage{geometry}                		% See geometry.pdf to learn the layout options. There are lots.
\geometry{letterpaper}                   		% ... or a4paper or a5paper or ... 
%\geometry{landscape}                		% Activate for for rotated page geometry
%\usepackage[parfill]{parskip}    		% Activate to begin paragraphs with an empty line rather than an indent
\usepackage{graphicx}				% Use pdf, png, jpg, or eps§ with pdflatex; use eps in DVI mode
								% TeX will automatically convert eps --> pdf in pdflatex		
\usepackage{amsmath, amsthm, amssymb, amsfonts}
\usepackage{mathabx}
\usepackage{makeidx}
\usepackage[symbol,perpage]{footmisc}
\usepackage[pdftex,bookmarks=true,pdfborder={0 0 0},colorlinks=true,linkcolor=blue,citecolor=blue]{hyperref}
\usepackage{setspace}

\newtheorem{definition}{Definition}[chapter]
\newtheorem{example}{Example}[chapter]
\newtheorem{remark}{Remark}[chapter]
\newtheorem{lemma}{Lemma}[chapter]
\newtheorem{corollary}{Corollary}[chapter]
\newtheorem{proposition}{Proposition}[chapter]
\newtheorem{axiom}{Axiom}[chapter]
\newtheorem{theorem}{Theorem}[chapter]
\newtheorem{problem}{Problem}[chapter]

\newcommand{\paren}[1]{\left(#1\right)}
\newcommand{\set}[1]{\left\{#1\right\}}
\newcommand{\var}[1]{\mathrm{Var}\left(#1\right)}
\newcommand{\cov}[1]{\mathrm{Cov}\left(#1\right)}

\makeindex 

\title{Probability Notes}
\author{Michael Conlen}
%\date{}							% Activate to display a given date or no date

\begin{document}
\frontmatter
\maketitle
\tableofcontents
\setcounter{tocdepth}{4}

\mainmatter
\begin{spacing}{1.618}
\chapter{Introduction to Probability}
	\section{Probability Spaces}
		\begin{definition}[Probability Space]\index{probability space}
			A probability space is a measure space such that the measure of the entire space is equal to one. We denote the space $(\Omega, \mathcal{F}, P)$ or alternatively $(S, E, P)$ where
			\begin{itemize}
				\item $\Omega$ is an arbitrary non-empty set called the \emph{sample space}\index{sample space}. 
				\item $\mathcal{F}$ is the $\sigma$-algebra called the \emph{event space}\index{event space} where elements of $\mathcal{F}$ are called \emph{events}\index{events}. 
				\item The measure $P:\mathcal{F} \to [0, 1]$ is called the \emph{probability}\index{probability}. 
			\end{itemize}
		\end{definition}
	
		\begin{remark}[Measure Space]\index{measure space}
			Since a probability space is a measure space we note that $F$ contains the entire space and is closed under complements and countable unions. The probability measure $P$ is countably additive (also called $sigma$-additive and the measure of the entire space is one. 
		\end{remark}

	\section{Axioms of Probability}	
		\begin{definition}[Axioms of Probability]\index{axioms of probability}
			There are three axioms of probability which hold true in any probability space
			\begin{axiom}
				For every event $A\subset \mathcal{F}$, $P(A) \geq 0$. 		
			\end{axiom} 
			\begin{axiom}
				For every sample space $\Omega$, $P(\Omega)=1$. 
			\end{axiom}
			\begin{axiom}
				For every infinite sequence of disjoint events $A_1, A_2, \dots$, 
				\begin{alignat}{4}
					P\left(\bigcup_{i=1}^\infty A_i\right)=\sum_{i=1}^\infty P(A_i)
				\end{alignat}
			\end{axiom}
		\end{definition}

	\section{Useful Corollaries}
	
		\begin{theorem}
			$P(\emptyset)=0$		
		\end{theorem}
		\begin{theorem}
			For every finite sequence of $n$ disjoint events $A_1, A_2, \dots, A_n$, 
				\begin{alignat}{4}
					P\left(\bigcup_{i=1}^n A_i\right)=\sum_{i=1}^n P(A_i)
				\end{alignat}
		\end{theorem}
		\begin{theorem}
			For every event $A$, $P(A^c)=1-P(A)$. 
		\end{theorem}
		
		\begin{theorem}
			If $A\subset B$ then $P(A)\leq P(B)$. 
		\end{theorem}
		
		\begin{theorem}
			For every event $A$, $0\leq A\leq 1$. 
		\end{theorem}
		
		\begin{theorem}
			For every two events $A$ and $B$, 
			\begin{alignat}{4}
				P(A\cap B^c)=P(A)-P(A\cap B)
			\end{alignat}
		\end{theorem}
		
		\begin{theorem}\label{Theorem:1:1}
			For every two events $A$ and $B$, 
			\begin{alignat}{4}
				P(A \cup B) = P(A)+P(B)-P(A\cap B)
			\end{alignat}
		\end{theorem}
		
		\begin{remark}
			Theorem \ref{Theorem:1:1} says that the probability of a union of events is equal to their sum minus their intersection. 
		\end{remark}
		
		\begin{theorem}[Bonferroni Inequality]\index{Bonferroni inequality}
			For rall events $A_1, A_2, \dots, A_n$, 
			\begin{alignat}{4}
				P(\left(\bigcup_{i=1}^n A_i\right) \leq \sum_{i=1}^nP(A_i)
			\end{alignat}
			and
			\begin{alignat}{4}
				P(\left(\bigcap{i=1}^n A_i\right) \geq 1-\sum_{i=1}^nP(A_i^c)
			\end{alignat}
		\end{theorem}
\chapter{Conditional Probabilities}
	\section{Definition of Conditional Probability}
	
		\begin{definition}[Conditional Probability]\index{conditional probability}
			Suppose we know that some event $B$ occurred and we wish to know the probability that some other event $A$ has occurred; further assume that $P(B)>0$, then we define the \emph{conditional probability} of $A$ given $B$ as 
			\begin{alignat}{4}
				P(A\mid B) = \frac{P(A\cap B)}{P(B)}
			\end{alignat}
			NB: If $P(B)=0$ then $P(A\mid B)$ is undefined. 
		\end{definition}
		
		\begin{theorem}
			Let $A$ and $B$ be events and $P(B)>0$, then
			\begin{alignat}{4}
				P(A\cap B)=P(B)P(A\mid B)
			\end{alignat}
			and if $P(A)>0$, then 
			\begin{alignat}{4}
				P(A\cap B)=P(A)P(B\mid A)
			\end{alignat}
		\end{theorem}
	\section{Partitions}
	

		\begin{definition}[Partition]\index{partition}
			Let $\Omega$ be a sample space and let $A_1, A_2, \dots, A_k$ be a set of $k$ events such that $A_i \cap A_j = \emptyset$ for all $i$, $j$ and that $\bigcup_{i=1}^k A_i = \Omega$; then we say that $\left\{A_i\right\}_{i=1}^k$ form a \emph{partition} of $\Omega$
		\end{definition}
		
		\begin{theorem}[Law of total probability]\index{law of total probability}
			Suppose that $\set{B_i}_{i=1}^k$ partition $\Omega$ and that $P(B_i)>0$ for all $0<i<k$; then for every event $A$
			\begin{alignat}{4}
				P(A)=\sum_{i=1}^kP(B_i)P(A\mid B_i)
			\end{alignat}
		\end{theorem}

	\section{Independent Events}
	
		\begin{definition}[Independent Event]\index{independent events}
			Two events, $A$ and $B$ are \emph{independent} if $P(A\cap B) = P(A)P(B)$. 
		\end{definition}
		
		\begin{theorem}
			If two events $A$ and $B$ are independent then $A$ and $B^c$ are also independent. 
		\end{theorem}

		

\printindex
\end{spacing}
\end{document}  